\documentclass[../../main]{subfiles}

\begin{document}

\section{Boas Ordens}
\label{sec:matematica:relacoes:boas-ordens}

\begin{defn}[Boa Ordem]
  \label{def:matematica:relacoes:boa-ordem}

  Seja ((A,\le)) um conjunto totalmente ordenado.

  Dizemos que (\le) é uma boa ordem quando todo subconjunto não vazio
  de (A) possui um menor elemento.

  Formalmente,

  [
    (\forall B\subseteq A)
    \Bigl(
      B\neq\varnothing
      \Rightarrow
      (\exists m\in B)
      (\forall x\in B)
      (m\le x)
    \Bigr).
  ]

\end{defn}

\begin{prop}
  \label{prop:matematica:relacoes:boa-ordem-menor-unico}

  Em uma boa ordem, todo subconjunto não vazio possui um único menor
  elemento.

\end{prop}

\begin{proof}

  A existência decorre da definição.

  A unicidade segue da antissimetria.

\end{proof}

\begin{prop}
  \label{prop:matematica:relacoes:boa-ordem-bem-fundada}

  Toda boa ordem é uma relação bem-fundada.
\end{prop}

\begin{proof}

  Seja (B\neq\varnothing).

  Como ((A,\le)) é uma boa ordem, existe um menor elemento
  (m\in B).

  Logo não existe elemento de (B) estritamente menor que (m).

  Portanto (m) é minimal.

\end{proof}

\begin{prop}[Princípio da Boa Ordem]
  \label{prop:matematica:relacoes:principio-boa-ordem}

  Todo subconjunto não vazio de uma boa ordem possui um menor elemento.

\end{prop}

\begin{proof}

  Trata-se da própria definição de boa ordem.

\end{proof}

\begin{prop}[Indução em Boas Ordens]
  \label{prop:matematica:relacoes:inducao-boa-ordem}

  Seja ((A,\le)) uma boa ordem.

  Seja (P(x)) uma proposição definida para todo (x\in A).

  Se

  [
    (\forall x\in A)
    \Bigl(
      (\forall y<x)
      P(y)
      \Rightarrow
      P(x)
    \Bigr),
  ]

  então

  [
    (\forall x\in A)
    P(x).
  ]

\end{prop}

\begin{proof}

  Suponha que exista um elemento para o qual (P) seja falsa.

  Considere

  [
    S
    =
    {
      x\in A
      \mid
      \neg P(x)
    }.
  ]

  Como (S\neq\varnothing), pela boa ordem existe um menor elemento
  (m\in S).

  Pela minimalidade de (m), todo elemento (y<m) satisfaz (P(y)).

  Pela hipótese,

  [
    P(m).
  ]

  Contradição.

  Logo,

  [
    (\forall x\in A)
    P(x).
  ]

\end{proof}

\begin{ex}

  O conjunto

  [
    \mathbb N
  ]

  munido da ordem usual

  [
    \le
  ]

  é uma boa ordem.

\end{ex}

\begin{obs}

  O princípio da indução matemática pode ser interpretado como um caso
  particular do princípio da indução sobre boas ordens.

\end{obs}

\end{document}
